\documentclass[11pt]{article}
\usepackage[margin=1in]{geometry}
\usepackage{amsmath,amssymb,amsthm,booktabs}
\newtheorem{theorem}{Theorem}
\newtheorem{corollary}{Corollary}
\newtheorem{proposition}{Proposition}
\newtheorem{remark}{Remark}
\title{Overlap-Correct Collar Charging and the Coherence--Charge Trade-off}
\author{}
\date{21 July 2026}
\begin{document}
\maketitle

\begin{abstract}
We derive the multi-link free-energy consequence of continuous collar failure
with overlap multiplicity retained.  Combining it with connected-witness
entropy exposes a quantitative incompatibility: inside the current nonlinear
window, the coherence radius small enough for the coherent conditional bound
is smaller than the radius needed to suppress incoherent-collar contours, even
under ideal overlap one.  The deficit yields explicit repair targets.
\end{abstract}

\section{Overlap-correct deterministic charge}

Let $W_1,\dots,W_k$ be selected collar-failure witnesses.  Stage 16 implies
that each witness has incident squared-curvature energy at least $\eta^2/2$.
Assume each plaquette occurs in at most $\nu$ selected incident families.

\begin{theorem}[Multi-witness charge]
The global squared-curvature energy and its heat-kernel Gibbs exponent satisfy
\[
 E_{\rm global}\ge \frac{k\eta^2}{2\nu},
 \qquad
 \frac{E_{\rm global}}{2t}\ge\frac{k\eta^2}{4\nu t}.
\]
\end{theorem}

\begin{proof}
Summing the $k$ local inequalities gives $k\eta^2/2$.  Every global
plaquette-energy term is counted at most $\nu$ times, so the sum of local
energies is at most $\nu E_{\rm global}$.  Division by $\nu$ and then $2t$
gives the result.
\end{proof}

The factor $\nu$ is necessary.  Repeating one identical witness $k$ times has
$\nu=k$ and produces only one local charge; a naive disjoint sum overstates
the exponent by exactly $k$.

\section{Connected entropy}

Let $D$ be the witness adjacency degree and let $e^{\alpha k}$ be the retained
state weight.  With
\[
 a_\nu(\eta)=\frac{\eta^2}{4\nu t},\qquad
 \theta=e^\alpha D^2e^{-a_\nu(\eta)},
\]
the standard rooted connected-set count gives, for minimum size $k_0$,
\[
 \sum_{k\ge k_0}e^{\alpha k}D^{2(k-1)}e^{-ka_\nu(\eta)}
 \le
 \frac{e^{\alpha k_0}e^{-k_0a_\nu(\eta)}D^{2(k_0-1)}}{1-\theta}
\]
whenever $\theta<1$.

\begin{corollary}[Entropy threshold]
Strict convergence requires
\[
 \eta>\eta_{\rm ent}(\nu)
 :=\sqrt{4\nu t\bigl(\alpha+2\log D\bigr)}.
\]
The executable module also solves the stronger threshold at which the negative
logarithm of the rooted sum exceeds $6.71292101139$.
\end{corollary}

\section{The common-parameter test}

Freeze $t=3\cdot10^{-6}$, $D=18$, $\alpha=0.1$, and the maximum sharp local
radius $0.011453785959974176$.  At twice that radius, compute the largest
$\eta_{\rm coh}$ whose optimized coherent conditional bound fits the remaining
$K+R$ budget.  Separately compute $\eta_{\rm FE}$ required by the rooted
incoherent contour with $k_0=2$.

\begin{proposition}[Quantified trade-off obstruction]
The accompanying interval-certified calibration gives
\[
 \eta_{\rm coh}<\eta_{\rm FE}
\]
already for ideal overlap multiplicity $\nu=1$.  Therefore the present
coherent estimate, deterministic collar charge, and degree-$18$ connected-set
count have no common coherence parameter inside the certified nonlinear
window.  Any $\nu>1$ increases the deficit.
\end{proposition}

This proposition is a no-go for the current proof mechanism, not for
Yang--Mills theory.  The calibration translates the deficit into two positive
targets: reduce the effective connected branching degree below the reported
ceiling, or amplify the deterministic exponent by the reported factor.  A
sector-aware polymer can also change the counting law rather than either local
constant.

\begin{remark}[Claim boundary]
The contour theorem is conditional on a witness selection with declared
overlap and on a normalized partition comparison.  The programme has not yet
constructed this selection uniformly over adaptive histories, boundary states,
volumes, orbit charts, and topological sectors.  Hence neither physical basin
entry nor a mass gap follows.
\end{remark}
\end{document}
